\exnumber{10}
\extitle{UI III - Ex 11 - Consistency}
\exfiles{ex10/Figma file, named Name\_FirstName\_DeliverableName\_Date\_VersionNumber}
\exauthorize{Those named in the exercise}
\makeheaderfiles

\textbf{Instructions:}

Pick a website or an app of your choice*. Highlight 5 signs of consistency across the pages.

*Airbnb, Google, Apple, Amazon, Facebook, Instagram and WhatsApp are not accepted.

\textbf{Deliverable:}

\begin{itemize}
\item A Figma file gathering at least 5 signs of consistency.
\item Upload the Figma file on Github, labeled as follows: "Name\_FirstName\_DeliverableName\_Date\_VersionNumber". For example, it can be labeled as "Doe\_John\_MidFiPrototype\_05242024\_V1".
\item Add a title within the file.
\item Pay attention to the global aspect of the file. It must be clear, simple and easy to read. You can get inspiration on canva.com but don't overload your design with too much details!
\end{itemize}

\textbf{Resources:}

\begin{itemize}
\item 7 tips for designing consistency \url{https://designshack.net/articles/graphics/7-tips-for-designing-consistency/}
\item The value of consistent design \url{https://www.invisionapp.com/inside-design/consistent-design/}
\item Consistency and Standards - Examples of Nielsens's Design Heuristics \url{https://medium.com/@gregoralbrecht/consistency-and-standards-nielsens-design-heuristic-explained-ac91f450fd8}
\item The importance of consistency in design work \url{https://yesimadesigner.com/the-importance-of-consistancy-in-design-work/}
\item Examples of signs of consistency: 
\begin{itemize}
\item colors
\item same shaped buttons
\item alignment
\item consistent product pages (with identical frameworks)
\item etc.
\end{itemize}
\end{itemize}

